\begin{ExerciseBlock}
\ExerciseHeading{06}{Bild- und Gegenstandsgröße}{geometrical_and_wave_optics_thin_biconvex_lens}
Eine dünne bikonvexe Glaslinse in Luft ($n=1{,}5$) bildet einen Gegenstand auf einen Schirm ab. Die Beträge der Krümmungsradien sind $R_1=\qty{30}{\centi\metre}$ und $R_2=\qty{50}{\centi\metre}$. Das Bild hat die halbe Größe des Gegenstands.

Wie weit ist die Linse vom Gegenstand und wie weit vom Schirm entfernt?
\end{ExerciseBlock}
\UseExerciseSolution{geometrical_and_wave_optics_thin_biconvex_lens}
